
Production RAG systems filter billions of web documents using GPT-2 perplexity, a quality metric inherited from language model pretraining that prioritizes narrative fluency and grammatical coherence. This practice assumes pretraining-optimal corpora are also retrieval-optimal. However, retrieval operates under different constraints: where pretraining values text a model can internalize during training, retrieval values text a model can locate and extract at inference time. A document with high perplexity due to technical jargon or structured tables may nonetheless contain the precise entities needed to answer factoid queries. If retrieval utility diverges from pretraining utility, current filtering strategies may systematically exclude retrieval-optimal documents.

Recent work on pretraining corpus quality demonstrates that model-based filtering outperforms heuristic approaches. DataComp-LM showed that training classifiers on high-quality pretraining examples enabled a 7B model to reach 64\% MMLU accuracy with 40\% less compute \citep{Li2024}. FineWeb-Edu demonstrated that educational quality classifiers produce dramatic gains on knowledge-intensive tasks---+5.0pp on MMLU, +4.5pp on ARC \citep{Penedo2024}. These methods optimize quality for language model pretraining, measuring success through next-token prediction loss and downstream task performance after training. Whether these quality signals transfer to retrieval---where success requires locating relevant information in a corpus at inference time---remains an open empirical question.

We investigate whether retrieval-quality signals diverge from pretraining quality by training classifiers on BEIR retrieval success examples. Our approach uses stratified sampling to enforce divergence: we oversample documents with low educational quality (high perplexity) but high BEIR retrieval relevance, forcing classifiers to learn signals orthogonal to pretraining fluency. We hypothesized that retrieval-optimal documents would exhibit higher factual entity density, operationalized through named entity recognition, and that this density would translate to measurable improvements specifically on semantic queries where lexical matching fails.

Our experiments yielded mixed results. In proof-of-concept validation on BEIR Natural Questions, retrieval-quality filtering achieved Recall@10 of 0.520 compared to the perplexity baseline's 0.470---a delta of +0.050 (+10.6\% relative improvement), exceeding our +0.03 gate threshold. This provides initial evidence that retrieval-specific filtering may be feasible. However, mechanism testing refuted our hypothesized causal pathway: retrieval-selected documents exhibited entity density ratio of 0.973 (2.7\% decrease vs. 15\% target increase), and showed no preferential advantage on semantic queries ($\Delta Recall$ = 0.00pp vs. $\geq 4pp$ target). The improvement appears real at proof-of-concept scale; the mechanism driving that improvement remains unidentified.

This work makes three contributions. First, we provide an initial controlled empirical test suggesting retrieval-quality signals may diverge from pretraining quality, demonstrated through proof-of-concept validation showing +10.6\% Recall@10 improvement over perplexity baselines. Second, we falsify the entity density hypothesis, showing that NER-based factual density does not correlate with BEIR retrieval quality (ratio 0.973 < 1.15 threshold), thereby narrowing the mechanistic search space. Third, we document methodological challenges in corpus-scale retrieval experiments, including corpus sampling issues that affect query coverage and the distinction between proof-of-concept validation and full-scale confirmation.

Our findings redirect research toward alternative mechanisms including query-document semantic alignment, answer-bearing sentence structures, or non-entity informativeness. The field needs theoretical frameworks for retrieval quality independent of pretraining paradigms.
